\section{Blind Signature Construction}
\label{sec:blind-signature}

This section presents the vSIS-based blind signature that underlies ARC-SPRUCE. It combines the Ajtai commitment of
Section~\ref{subsec:commitments}, the trapdoor rational hash of Definition~\ref{def:rational-hash}, and the NIZK-ARK of
Definition~\ref{def:nizk-ark}. The signed message is always
\[
  \vecm=\vecmu\Vert \veck,
\]
where \(\vecmu\) denotes the attribute payload and \(\veck\) is a holder-generated secret PRF key used later for rate
limiting.

\subsubsection{Randomizing the rational-hash witness}
In the protocol, \(\Holder\) first commits to
\[
  \vecw=[\vecm\Vert \vecr_0^{\Holder}\Vert \vecr_1^{\Holder}\Vert \bar\vecr]^\top,
\]
where \((\vecr_0^{\Holder},\vecr_1^{\Holder})\) are the holder contributions to the rational-hash witness and
\(\bar\vecr\) is fresh commitment randomness. To satisfy the commitment bounds, we require
\(\|\vecw\|_\infty\le \BoundMsg\).

The issuance process then uses a two-party randomness generation step: the issuer samples
\((\vecr_0^{\Issuer},\vecr_1^{\Issuer})\), and the joint witness is obtained coefficient-wise by
\[
  \vecr_0=\centerfunc(\vecr_0^{\Holder}+\vecr_0^{\Issuer}),\qquad
  \vecr_1=\centerfunc(\vecr_1^{\Holder}+\vecr_1^{\Issuer}).
\]
Appendix~\ref{app:defs} records the centering map and its uniformity property.

\subsubsection{Blind signature protocol}
\label{subsec:issuance-protocol}
We assume \(\nizkscheme=(\ZKProve,\ZKVerify)\) is a NIZK for the relations below,
\(\ComScheme=(\ComSetup,\Commit,\ComVerify)\) is the Ajtai commitment scheme, and
\[
  \TDRH=(\Setup,\TrapGen,\mathsf{Eval},\SampPre)
\]
is the trapdoor rational-hash interface from Definition~\ref{def:bbs-rational-hash}. Public setup outputs the hash key
\(B\), while issuer key generation outputs the trapdoored signing matrix \(\mA\).

\paragraph{Abstract and compiled issuance relations.}
We distinguish the ideal issuance relation from the relation actually passed to the pre-sign proof system.

We write \(R_{\mathrm{Issue}}^{\mathrm{abs}}(x,w)=1\) iff
\[
\begin{aligned}
  x&=\bigl(\ComPublicParam,\vecc,T,(\vecr_0^{\Issuer},\vecr_1^{\Issuer}),B\bigr),\\
  w&=(\vecm,\vecr_0^{\Holder},\vecr_1^{\Holder},\bar\vecr),\\
  \ComPublicParam&=(\mACom,\BoundMsg),\\
  \vecc&=\mACom\cdot[\vecm\Vert \vecr_0^{\Holder}\Vert \vecr_1^{\Holder}\Vert \bar\vecr]^\top,\\
  \|\vecm\|_\infty,\|\vecr_0^{\Holder}\|_\infty,\|\vecr_1^{\Holder}\|_\infty,\|\bar\vecr\|_\infty&\le \BoundMsg,\\
  \vecr_0&=\centerfunc(\vecr_0^{\Holder}+\vecr_0^{\Issuer}),\\
  \vecr_1&=\centerfunc(\vecr_1^{\Holder}+\vecr_1^{\Issuer}),\\
  T&=h_{\vecm,(\vecr_0,\vecr_1)}(B).
\end{aligned}
\]

We write \(R_{\mathrm{Issue}}^{\mathrm{cmp}}(x,w)=1\) iff the same conditions hold, except that the last line is replaced by
\[
  (B_3-\vecr_1)\Had T-\bigl(B_0+B_1\cdot \vecm+B_2\cdot \vecr_0\bigr)=0.
\]

The issued signature relation remains the abstract relation
\[
R_{\mathrm{Sign}}(x,w)=1
\iff
x=(\vecm,\mA,B,\BoundMsg,\BoundSig),\quad w=(\vecu,\vecr_0,\vecr_1),
\]
\[
  \|\vecr_0\|_\infty,\|\vecr_1\|_\infty \le \BoundMsg,
  \qquad
  \|\vecu\|_\infty \le \BoundSig,
\]
\[
  h_{\vecm,(\vecr_0,\vecr_1)}(B)\neq\bot,
  \qquad
  \mA\vecu=h_{\vecm,(\vecr_0,\vecr_1)}(B).
\]

\begin{lemma}\label{lem:issue-abs-cmp}
If an honest holder first checks that \(B_3-\vecr_1\) is invertible in every CRT slot and then computes
\(T=h_{\vecm,(\vecr_0,\vecr_1)}(B)\), every witness satisfying \(R_{\mathrm{Issue}}^{\mathrm{abs}}\) also satisfies
\(R_{\mathrm{Issue}}^{\mathrm{cmp}}\). The converse need not hold.
\end{lemma}

\begin{proof}
This is immediate from Lemma~\ref{lem:hash-cleared-admissible} and Remark~\ref{rem:hash-cleared-not-sufficient}.
\end{proof}

We now define the blind-signature algorithms.
\begin{description}
  \item[$\Setup(\SecParam)$:] Fix the global ring parameters. Run
  \[
    (\mACom,\BoundMsg)\gets \ComScheme.\ComSetup(\SecParam)
  \]
  and
  \[
    (B,\mathcal{M},\mathcal{X},\delta,\BoundMsg,\BoundSig)\gets \TDRH.\Setup(\SecParam).
  \]
  Output
  \[
    \PublicParamsBS=(\mACom,B,\mathcal{M},\mathcal{X},\delta,\BoundMsg,\BoundSig).
  \]

  \item[$\IKeyGen()$:] Compute \((\mA,\trapdoor)\gets \TDRH.\TrapGen(\SecParam)\) and output
  \[
    \vk=\mA,\qquad \sk=\trapdoor.
  \]

  \item[$\Issue$:] The issuer \(\Issuer\) holds \(\sk\); the holder \(\Holder\) holds \(\vecmu,\vk\).
  \begin{enumerate}[leftmargin=1.25em]
    \item \(\Holder\) samples \(\veck,\vecr_0^{\Holder},\vecr_1^{\Holder},\bar\vecr\) coefficient-wise in
    \([ -\BoundMsg,\BoundMsg]\) and sets \(\vecm=\vecmu\Vert \veck\).
    \item \(\Holder\) computes
    \[
      \vecc\gets \mACom\cdot[\vecm\Vert \vecr_0^{\Holder}\Vert \vecr_1^{\Holder}\Vert \bar\vecr]^\top
    \]
    and sends \(\vecc\) to \(\Issuer\).
    \item \(\Issuer\) samples \(\vecr_0^{\Issuer},\vecr_1^{\Issuer}\) coefficient-wise in
    \([ -\BoundMsg,\BoundMsg]\) and sends them to \(\Holder\).
    \item \(\Holder\) computes
    \[
      \vecr_0\gets \centerfunc(\vecr_0^{\Holder}+\vecr_0^{\Issuer}),\qquad
      \vecr_1\gets \centerfunc(\vecr_1^{\Holder}+\vecr_1^{\Issuer}),
    \]
    checks whether \(h_{\vecm,(\vecr_0,\vecr_1)}(B)=\bot\), and aborts/restarts if so. Otherwise it computes the public target
    \[
      T\gets h_{\vecm,(\vecr_0,\vecr_1)}(B).
    \]
    \item \(\Holder\) computes \(\nizkIssue\gets \nizkscheme.\ZKProve(w,x)\) for the compiled relation
    \(R_{\mathrm{Issue}}^{\mathrm{cmp}}\), with witness
    \[
      w=(\vecm,\vecr_0^{\Holder},\vecr_1^{\Holder},\bar\vecr)
    \]
    and statement
    \[
      x=(\ComPublicParam,\vecc,T,(\vecr_0^{\Issuer},\vecr_1^{\Issuer}),B),
    \]
    and sends \((T,\nizkIssue)\) to \(\Issuer\).
    \item \(\Issuer\) verifies \(\nizkIssue\) for \(R_{\mathrm{Issue}}^{\mathrm{cmp}}\). If verification fails, it aborts.
    Otherwise it computes
    \[
      \vecu\gets \TDRH.\SampPre(\trapdoor,T)
    \]
    and sends \(\vecu\) to \(\Holder\).
    \item \(\Holder\) checks \(\|\vecu\|_\infty\le \BoundSig\) and \(\mA\vecu=T\). If so, it outputs
    \[
      \sig=(\vecu,\vecr_0,\vecr_1)
    \]
    on message \(\vecm=\vecmu\Vert \veck\); otherwise it aborts.
  \end{enumerate}

  \item[$\Verify(\vk,\vecm,\sig)$:] Let \(\vk=\mA\). Parse \(\sig=(\vecu,\vecr_0,\vecr_1)\). Accept iff
  \[
    \bigl((\vecm,\mA,B,\BoundMsg,\BoundSig),(\vecu,\vecr_0,\vecr_1)\bigr)\in R_{\mathrm{Sign}}.
  \]
\end{description}

For ARC, the holder does not reveal \(\sig\) directly. Section~\ref{sec:arc-construction} instead proves the compiled
post-sign relation of Section~\ref{sec:smallwood-model} in zero knowledge while keeping \(\veck,\vecr_0,\vecr_1\) hidden
and binding \(\veck\) to the PRF tag.

\subsubsection{Security}
\begin{proposition}[Correctness]
Assume the pre-sign NIZK is complete and the trapdoor sampler is correct. Then an honest execution of \(\Issue\)
outputs a signature \(\sig\) that is accepted by \(\Verify\), except with negligible failure probability coming from the
trapdoor sampler or an inadmissible hash evaluation.
\end{proposition}

\begin{proof}
By honest target derivation, the holder computes \(T\) only after checking that
\(h_{\vecm,(\vecr_0,\vecr_1)}(B)\neq\bot\), and then sets
\[
T=h_{\vecm,(\vecr_0,\vecr_1)}(B).
\]
Lemma~\ref{lem:issue-abs-cmp} implies that the honest witness satisfies the compiled relation
\(R_{\mathrm{Issue}}^{\mathrm{cmp}}\), so completeness of the pre-sign NIZK implies that the issuer accepts. By
correctness of \(\SampPre\), the returned \(\vecu\) satisfies \(\mA\vecu=T\) and the public shortness bound except with
negligible probability. Hence \(\sig=(\vecu,\vecr_0,\vecr_1)\) satisfies \(R_{\mathrm{Sign}}\).
\end{proof}

\begin{remark}[Security status of the blind protocol]\label{rem:blind-security-roadmap}
The non-blind hash-and-preimage signature layer underlying this protocol is justified by the DKLW25 reduction summarized
in Lemma~\ref{lem:vsis-unforgeability}. For the blind protocol itself, the only theorem-backed property proved in the
present draft is honest-party correctness. A proof of one-more unforgeability would still have to combine the
extractability of the pre-sign proof, the hiding of the Ajtai commitment, the algebraic correctness of the centering
step, and an interactive one-more hardness statement that genuinely matches the issuance oracle. A proof of blindness
would still have to combine the same commitment-hiding and proof zero-knowledge ingredients with a simulation argument
for the issuer view. Accordingly, the blind protocol is presented here as a construction with defined goals and an
explicit reduction roadmap, not as a layer for which blindness or one-more unforgeability has already been reduced.
\end{remark}
